%% Machine Learning and Agentic AI for Macroeconomic Research
%% Lecture 6: Heterogeneous Agent Models via Reinforcement Learning
%% Author: Zhigang Feng
%% Sources: 11_HA_model/Lec11_ML_HA_models.tex
%%          15-MFG/MFG.tex

\documentclass[10pt,english,aspectratio=169]{beamer}

\usepackage{ae,aecompl}
\usepackage[T1]{fontenc}
\usepackage{textcomp}
\usepackage[utf8]{inputenc}
\usepackage{babel}

\setcounter{secnumdepth}{3}
\setcounter{tocdepth}{3}

\usepackage{amsmath, amssymb, graphicx}
\usepackage{booktabs, multirow}
\usepackage{tcolorbox}
\tcbuselibrary{skins, breakable, theorems}
\usepackage{listings}
\usepackage{xcolor}

\lstset{
  basicstyle=\ttfamily\small,
  keywordstyle=\color{blue!80!black}\bfseries,
  commentstyle=\color{green!50!black}\itshape,
  stringstyle=\color{orange!80!black},
  backgroundcolor=\color{gray!8},
  frame=single, rulecolor=\color{gray!40},
  breaklines=true, showstringspaces=false,
  language=Python,
}

\ifx\hypersetup\undefined
  \AtBeginDocument{%
    \hypersetup{bookmarks=false,breaklinks=true,pdfborder={0 0 0},
      colorlinks=true,linkcolor=DarkText,urlcolor=RedTitle}%
  }
\else
  \hypersetup{bookmarks=false,breaklinks=true,pdfborder={0 0 0},
    colorlinks=true,linkcolor=DarkText,urlcolor=RedTitle}%
\fi

\makeatletter
\newcommand\makebeamertitle{\frame{\maketitle}}%
\AtBeginDocument{%
  \let\origtableofcontents=\tableofcontents
  \def\tableofcontents{\@ifnextchar[{\origtableofcontents}{\gobbletableofcontents}}
  \def\gobbletableofcontents#1{\origtableofcontents}%
}

\usetheme{metropolis}
\usefonttheme{professionalfonts}

\definecolor{LightGray}{RGB}{230,230,230}
\definecolor{RedTitle}{RGB}{0,0,200}
\definecolor{VeryLightGray}{RGB}{245,245,245}
\definecolor{DarkText}{RGB}{0,0,0}

\setbeamercolor{background canvas}{bg=white}
\setbeamercolor{title}{fg=RedTitle}
\setbeamercolor{author}{fg=DarkText}
\setbeamercolor{date}{fg=DarkText}
\setbeamercolor{frametitle}{bg=VeryLightGray,fg=DarkText}
\setbeamercolor{normal text}{fg=DarkText}

\setbeamersize{text margin left=5mm, text margin right=5mm}
\addtobeamertemplate{frametitle}{}{\vspace{-0.5em}}
\newcommand{\reducevspace}{\vspace{-0.5cm}}

\setbeamertemplate{navigation symbols}{}
\setbeamertemplate{footline}{%
  \hfill%
  \usebeamercolor[fg]{page number in head/foot}%
  \usebeamerfont{page number in head/foot}%
  \insertframenumber\,/\,\inserttotalframenumber\kern1em\vskip2pt%
}

\makeatother

\graphicspath{{../../11_HA_model/}{../../15-MFG/}}

\title[ML \& Agentic AI for Macro Research]{Machine Learning and Agentic AI for Macroeconomic Research\\
6: Heterogeneous Agent Models via Reinforcement Learning\\[0.5em]}
\author{Zhigang Feng}
\date{Spring 2026}

\begin{document}

\makebeamertitle

%=============================================================================
% AGENDA
%=============================================================================

\begin{frame}{Lecture Agenda}
\begin{enumerate}
  \item \textbf{Aiyagari--Bewley--Huggett} --- incomplete-markets HA model and stationary equilibrium
  \medskip
  \item \textbf{Computation of stationary equilibrium} --- the classical three-step algorithm
  \medskip
  \item \textbf{Why ML?} --- motivating example on optimal growth, neural-network approximation
  \medskip
  \item \textbf{Euler-equation methods} --- residual losses, supervised learning, actor--critic
  \medskip
  \item \textbf{Global solution via DNN} --- distribution as state, value/policy networks
  \medskip
  \item \textbf{Krusell--Smith (1998)} --- aggregate uncertainty and bounded-rationality moments
  \medskip
  \item \textbf{DeepHAM and DEQN} --- modern deep-learning solutions for HA models
  \medskip
  \item \textbf{OLG with aggregate uncertainty} --- life-cycle structure, homotopy training
  \medskip
  \item \textbf{Mean-Field Games} --- continuous-time HA, HJB \& KFE, link to RL
  \medskip
  \item \textbf{Lab preview} --- Krusell--Smith via Deep Equilibrium Nets in PyTorch
\end{enumerate}
\end{frame}

%=============================================================================
% PART I: AIYAGARI / BEWLEY / HUGGETT --- SETUP & EQUILIBRIUM
%=============================================================================

\section{Aiyagari--Bewley--Huggett Setup}

\begin{frame}{Why heterogeneous agents?}
\begin{itemize}
\item The representative-agent paradigm assumes \emph{complete markets} and a single decision maker.
\medskip
\item Many empirically important questions require \textbf{heterogeneity}:
  \begin{itemize}
  \item Wealth and income inequality
  \item Precautionary savings and the demand for insurance
  \item Distributional consequences of monetary and fiscal policy
  \item Heterogeneous responses to aggregate shocks
  \end{itemize}
\medskip
\item HA models combine:
  \begin{itemize}
  \item Idiosyncratic income (or productivity, mortality, $\ldots$) shocks
  \item Incomplete asset markets (a single risk-free bond, borrowing constraint)
  \item General equilibrium in factor prices
  \end{itemize}
\medskip
\item This lecture: Aiyagari $\to$ Krusell--Smith $\to$ deep-learning solutions $\to$ Mean-Field Games.
\end{itemize}
\end{frame}

\begin{frame}{Aiyagari (1994, QJE): heterogeneity among households}
\begin{itemize}
\item A continuum of \emph{ex-ante identical}, infinitely-lived households.
\medskip
\item Each receives an idiosyncratic, uninsurable labour-productivity shock $\epsilon_t$ following a Markov process.
\medskip
\item Households self-insure by accumulating a single risk-free asset $a_t$.
\medskip
\item A natural or ad-hoc \textbf{borrowing constraint} $a_t \geq -\underline{a}$ keeps the asset distribution well-defined.
\medskip
\item Despite ex-ante symmetry, the realized history of shocks generates an \textbf{ex-post} distribution of wealth.
\medskip
\item Bewley (1980s, unpublished), Huggett (1993), Aiyagari (1994): the workhorse incomplete-markets HA framework.
\end{itemize}
\end{frame}

\begin{frame}{General equilibrium with ex-post heterogeneity}
\begin{itemize}
\item \textbf{Idiosyncratic shocks}: $\epsilon_t \in \mathcal{E}$, transition $\pi(\epsilon,\epsilon')$, with $\sum_{\epsilon}\epsilon\,\pi^*(\epsilon)=L$ constant.
\item \textbf{No aggregate shocks} in Aiyagari (we add them with Krusell--Smith later).
\item \textbf{Markets}:
  \begin{itemize}
  \item Risk-free bond (or capital) at price $1/(1+r)$
  \item Labour supplied inelastically; wage $w$
  \end{itemize}
\item \textbf{Stationarity}: prices $(r,w)$ and the cross-sectional distribution $\lambda(a,\epsilon)$ are constant in time.
\item \textbf{Individual state}: $(a,\epsilon) \in \mathbb{R}_+ \times \mathcal{E}$.
\item \textbf{Aggregate state}: distribution $\lambda$ over individual states (in stationary equilibrium, $\lambda=\lambda^*$).
\end{itemize}
\end{frame}

\begin{frame}{Household's recursive problem}
For given prices $(r,w)$ and stationary distribution $\lambda$, each household solves
\[
v(a,\epsilon;\lambda) \;=\; \max_{c,\,a'}\!\left\{u(c) + \beta \sum_{\epsilon'} v(a',\epsilon';\lambda)\,\pi(\epsilon,\epsilon')\right\}
\]
subject to
\[
c + a' = (1+r)\,a + w\,\epsilon, \qquad a' \geq -\underline{a}, \qquad c \geq 0.
\]
\medskip
\textbf{Objects to solve for}
\begin{itemize}
\item Value function $v(a,\epsilon)$
\item Asset policy $a'(a,\epsilon)$
\item Consumption policy $c(a,\epsilon) = (1+r)a + w\epsilon - a'(a,\epsilon)$
\end{itemize}
\medskip
With CRRA utility $u(c)=c^{1-\sigma}/(1-\sigma)$ and standard assumptions on $\pi$, the value function is increasing, concave in $a$, and the policy is monotone.
\end{frame}

\begin{frame}{Budget constraint and asset structure}
\begin{itemize}
\item $a$ denotes a one-period risk-free claim that pays $(1+r)a$ next period.
\medskip
\item \textbf{Borrowing constraint}: $a' \geq -\underline{a}$.
  \begin{itemize}
  \item \emph{Natural} bound: $\underline{a} = w\epsilon_{\min}/r$ (largest debt repayable from minimum income).
  \item \emph{Ad-hoc} bound: any $\underline{a} \leq$ natural bound.
  \end{itemize}
\medskip
\item Without uncertainty, the household would set $\beta(1+r)=1$ and consumption-smooth perfectly. With uninsurable risk and a binding constraint, the \textbf{precautionary motive} drives extra saving.
\medskip
\item This precautionary saving raises aggregate capital relative to the complete-markets benchmark, and depresses $r$ below $\rho = 1/\beta - 1$.
\end{itemize}
\end{frame}

\begin{frame}{Firms}
A representative firm operates a constant-returns technology
\[
Y = F(K,L) = A K^{\alpha} L^{1-\alpha}.
\]
Competition in factor markets gives
\[
r = F_K(K,L) - \delta = \alpha A (K/L)^{\alpha-1} - \delta, \qquad
w = F_L(K,L) = (1-\alpha) A (K/L)^{\alpha}.
\]
\medskip
\begin{itemize}
\item Aggregate labour supply $L = \sum_\epsilon \epsilon\,\pi^*(\epsilon)$ is exogenous.
\item Aggregate capital $K$ is determined endogenously by household saving.
\item In equilibrium, $K$ and $r$ are jointly pinned down by household and firm decisions.
\end{itemize}
\end{frame}

\begin{frame}{Stationary recursive competitive equilibrium (SRCE) --- prices}
A \textbf{stationary recursive competitive equilibrium} consists of:
\begin{enumerate}
\item Factor prices $(r^*, w^*)$ derived from the firm's first-order conditions:
\[
r^* = \alpha A (K^*/L)^{\alpha-1} - \delta, \qquad w^* = (1-\alpha) A (K^*/L)^{\alpha}.
\]
\medskip
\item Value function $v(a,\epsilon)$ and policies $(c(a,\epsilon), a'(a,\epsilon))$ that solve the household's problem given $(r^*,w^*)$.
\medskip
\item A stationary cross-sectional distribution $\lambda^*(a,\epsilon)$.
\medskip
\item Market clearing conditions (next slide).
\end{enumerate}
\end{frame}

\begin{frame}{SRCE --- HH problem, market clearing, stationary distribution}
\textbf{Labour market clearing:}
\[
L = \sum_{\epsilon} \int \epsilon \, d\lambda^*(a,\epsilon).
\]

\textbf{Capital market clearing:}
\[
K^* = \sum_{\epsilon} \int a'(a,\epsilon) \, d\lambda^*(a,\epsilon).
\]

\textbf{Stationarity of the distribution:}
\[
\lambda^*(A',\epsilon') = \sum_{\epsilon} \pi(\epsilon,\epsilon') \int_{\{a:\,a'(a,\epsilon)\in A'\}} d\lambda^*(a,\epsilon).
\]
\medskip
That is, $\lambda^*$ is invariant under the Markov operator induced by the policy $a'(\cdot)$ and the shock transition $\pi$.
\end{frame}

\begin{frame}{Stationarity and the invariant measure}
\begin{itemize}
\item Time-invariant prices require a time-invariant cross-sectional distribution.
\medskip
\item The pair $(a'(\cdot), \pi)$ defines a Markov \textbf{transition operator} $Q$ on the state space $S = \mathbb{R}_+ \times \mathcal{E}$:
\[
(Qf)(a,\epsilon) = \sum_{\epsilon'} \pi(\epsilon,\epsilon') f(a'(a,\epsilon),\epsilon').
\]
\medskip
\item The invariant distribution $\lambda^*$ satisfies $\lambda^* = Q^* \lambda^*$ (where $Q^*$ is the adjoint).
\medskip
\item Existence and uniqueness of $\lambda^*$ are non-trivial: they depend on \textbf{Feller} continuity, \textbf{mixing}, and \textbf{monotonicity} of $Q$ --- the next few slides.
\end{itemize}
\end{frame}

\begin{frame}{Two markets}
\begin{itemize}
\item Aiyagari's equilibrium reduces to clearing in \textbf{two markets}:
  \begin{itemize}
  \item \emph{Labour market}: trivially satisfied because labour supply is inelastic.
  \item \emph{Capital market}: aggregate household savings = aggregate firm demand for capital.
  \end{itemize}
\medskip
\item All the action is in the capital market.
\medskip
\item Strategy: parametrise everything by $r$.
  \begin{itemize}
  \item For each candidate $r$, solve the household problem to get $a'(a,\epsilon;r)$.
  \item Compute the implied stationary distribution $\lambda_r$.
  \item Compute the supply of capital $A(r) = \int a\,d\lambda_r$.
  \item Compare to the demand $K(r)$ from the firm side; iterate.
  \end{itemize}
\end{itemize}
\end{frame}

\begin{frame}{Determination of $r$: demand for capital}
The firm's first-order condition gives a downward-sloping demand for capital:
\[
K(r) \;=\; L \left( \frac{r + \delta}{\alpha A} \right)^{\frac{1}{\alpha-1}}.
\]
\medskip
\begin{itemize}
\item As $r \uparrow$, capital becomes more expensive and demand falls.
\item $K(r)$ is continuous and strictly decreasing in $r$.
\item $K(r) \to \infty$ as $r \to -\delta$ and $K(r) \to 0$ as $r \to \infty$.
\end{itemize}
\medskip
This is the standard neoclassical capital demand schedule.
\end{frame}

\begin{frame}{Determination of $r$: supply of capital}
The household supply of capital is the cross-sectional average asset holding under the stationary distribution induced by $r$:
\[
A(r) \;=\; \sum_{\epsilon} \int a \, d\lambda_r(a,\epsilon).
\]
\medskip
\begin{itemize}
\item $A(r)$ is upward-sloping in $r$ (broadly).
\item In the complete-markets benchmark $\beta(1+r)=1$, $A(r)$ has a vertical asymptote at $r = 1/\beta-1=\rho$.
\item With incomplete markets, precautionary saving makes $A(r)$ finite \emph{below} $\rho$, so the equilibrium $r^* < \rho$.
\end{itemize}
\end{frame}

\begin{frame}{Equilibrium picture}
\begin{center}
\includegraphics[width=0.65\textwidth]{equilibrium}
\end{center}
\vspace{-0.5em}
\begin{itemize}
\item $K(r)$: firm demand (downward sloping).
\item $A(r)$: household supply (upward sloping, asymptotic to $\rho$ in the complete-markets case).
\item Equilibrium $r^*$ satisfies $K(r^*) = A(r^*)$, with $r^* < \rho$.
\end{itemize}
\end{frame}

\begin{frame}{Probability measures and the Markov operator $Q$}
\begin{itemize}
\item Let $\Lambda$ denote the space of probability measures on the Borel sets of $S = \mathbb{R}_+ \times \mathcal{E}$.
\medskip
\item Define the Markov operator $Q : \Lambda \to \Lambda$ by
\[
(Q\lambda)(B) = \int \mathbb{1}\{(a'(a,\epsilon),\epsilon') \in B\}\,\pi(\epsilon,d\epsilon')\,d\lambda(a,\epsilon).
\]
\medskip
\item Three properties of $Q$ are essential to study the equilibrium:
  \begin{itemize}
  \item \textbf{Feller property} $\to$ existence of an invariant measure.
  \item \textbf{Mixing property} $\to$ uniqueness.
  \item \textbf{Monotonicity} $\to$ convergence.
  \end{itemize}
\medskip
\item These are standard conditions in the Stokey--Lucas--Prescott (1989) framework.
\end{itemize}
\end{frame}

\begin{frame}{Feller property $\Rightarrow$ existence}
\textbf{Definition.} $Q$ has the \emph{Feller property} if it maps continuous bounded functions to continuous bounded functions:
\[
f \in C_b(S) \;\Rightarrow\; Q f \in C_b(S).
\]
\medskip
\begin{itemize}
\item Equivalently, the policy $a'(\cdot)$ is continuous in the state.
\medskip
\item Combined with compactness of the state space (or tightness of the orbit), Feller continuity guarantees the \textbf{existence of an invariant measure} via the Krylov--Bogolyubov theorem.
\medskip
\item In Aiyagari, the asset grid is bounded above by $w \epsilon_{\max}/r$ (since saving above this generates negative consumption probability), so $S$ is effectively compact.
\end{itemize}
\end{frame}

\begin{frame}{Mixing property $\Rightarrow$ uniqueness}
\textbf{Mixing condition.} For every two states $s_1, s_2 \in S$, there exist $n_1, n_2 \in \mathbb{N}$ such that the $n$-step distributions $Q^{n_1}\delta_{s_1}$ and $Q^{n_2}\delta_{s_2}$ have overlapping support.
\medskip
\begin{itemize}
\item Intuitively: starting from \emph{any} two initial states, the chain visits a common region in finite time.
\medskip
\item This rules out absorbing sets that would split the state space and admit multiple invariant measures.
\medskip
\item \textbf{Counterexample}: a deterministic policy that pushes any household with $a > a^*$ to $a \uparrow$ and any household with $a < a^*$ to $a \downarrow$ creates two disjoint absorbing sets and \emph{multiple} invariant distributions.
\medskip
\item In standard Aiyagari, idiosyncratic income shocks ensure mixing.
\end{itemize}
\end{frame}

\begin{frame}{Monotonicity $\Rightarrow$ convergence}
\textbf{Monotonicity.} If $\lambda \succeq \lambda'$ in the first-order stochastic dominance sense, then $Q\lambda \succeq Q\lambda'$.
\medskip
\begin{itemize}
\item Combined with Feller and mixing, monotonicity implies that iterating $Q$ from any initial distribution converges to the unique invariant measure $\lambda^*$.
\medskip
\item This is the foundation for \emph{simulating} the model: start from an arbitrary distribution, simulate forward many periods, and read off $\lambda^*$.
\medskip
\item \textbf{Counterexample}: a non-monotone policy can produce stationary distributions that are sensitive to the initial condition; iteration may cycle rather than converge.
\medskip
\item Concavity of the value function and convexity of the budget set yield monotone policies in Aiyagari.
\end{itemize}
\end{frame}

%=============================================================================
% PART II: COMPUTATION OF STATIONARY EQUILIBRIUM
%=============================================================================

\section{Computation of Stationary Equilibrium}

\begin{frame}{Aggregate precautionary savings, applications}
\begin{itemize}
\item Because precautionary saving raises $A(r)$, the equilibrium interest rate satisfies $r^* < \rho = 1/\beta-1$.
\medskip
\item Aggregate capital is \textbf{higher} than in the complete-markets benchmark.
\medskip
\item Standard applications of Aiyagari-style models:
  \begin{itemize}
  \item Wealth and income inequality (top 10\%, Gini, $\ldots$)
  \item Welfare effects of social insurance and tax/transfer policy
  \item Effects of inflation on wealth distribution
  \item Idiosyncratic-risk asset pricing
  \item Quantitative theories of entrepreneurship
  \end{itemize}
\medskip
\item Krueger--Mitman--Perri (2016), Kaplan--Moll--Violante (2018) HANK, $\ldots$
\end{itemize}
\end{frame}

\begin{frame}{Computation of SRCE: three-step algorithm}
\textbf{Goal:} solve for $(r^*, w^*, \lambda^*, v, a', c)$ jointly.
\medskip

\begin{enumerate}
\item \textbf{Step 1.} Fix a candidate $r$. Solve the household dynamic programming problem to obtain $v_r$ and policies $(a'_r, c_r)$.
\medskip
\item \textbf{Step 2.} Given $a'_r$, compute the stationary cross-sectional distribution $\lambda_r$ as the eigenvector of the Markov transition matrix.
\medskip
\item \textbf{Step 3.} Evaluate the excess demand for capital
\[
d(r) \;=\; K(r) - \int a\,d\lambda_r,
\]
and update $r$ via bisection (or Brent) until $|d(r)|<\text{tol}$.
\end{enumerate}
\medskip
This is the \textbf{nested fixed point} that has been the workhorse of HA computation since Aiyagari (1994).
\end{frame}

\begin{frame}{Step 1: solve the household problem for fixed $r$}
For given $(r,w)$, iterate the Bellman operator
\[
(T_r v)(a,\epsilon) \;=\; \max_{a' \geq -\underline{a}}\!\left\{u\big((1+r)a + w\epsilon - a'\big) + \beta \sum_{\epsilon'} v(a',\epsilon')\,\pi(\epsilon,\epsilon')\right\}
\]
on a grid for $a$.
\medskip
\begin{itemize}
\item Standard tools: VFI with monotonicity exploitation, Howard policy iteration, endogenous-grid method (Carroll), Coleman--Reffett operator.
\medskip
\item Concavity of $u$ and the convex budget set guarantee a unique optimum.
\medskip
\item Output: $a'_r(a,\epsilon)$, $c_r(a,\epsilon)$, $v_r(a,\epsilon)$.
\end{itemize}
\end{frame}

\begin{frame}{Step 2: stationary distribution from the transition matrix}
Discretise $a$ on a grid; combined with $\epsilon$, this gives a finite state space of size $N = N_a \cdot N_\epsilon$.
\medskip

The policy and shock transition induce a Markov matrix $\mathbf{Q}_r \in \mathbb{R}^{N\times N}$ with
\[
\mathbf{Q}_r\big[(a_i,\epsilon_j) \to (a_k,\epsilon_\ell)\big]
= \pi(\epsilon_j,\epsilon_\ell)\cdot\Pr\!\big(a'_r(a_i,\epsilon_j) = a_k\big).
\]
\medskip
The stationary distribution solves
\[
\lambda_r = \mathbf{Q}_r^\top \lambda_r, \qquad \mathbf{1}^\top \lambda_r = 1.
\]
\medskip
Computation: dominant left eigenvector of $\mathbf{Q}_r$, or simply iterate $\lambda \leftarrow \mathbf{Q}_r^\top \lambda$ until convergence (the monotonicity result above guarantees this works).
\end{frame}

\begin{frame}{Step 3: market clearing on $r$}
With $\lambda_r$ in hand, compute the household supply of capital
\[
A(r) \;=\; \sum_{i,j} a_i \,\lambda_r(a_i,\epsilon_j),
\]
and compare to firm demand:
\[
d(r) \;=\; K(r) - A(r).
\]
\medskip
\textbf{Update rule}:
\begin{itemize}
\item If $d(r) > 0$: capital is scarce, raise $r$.
\item If $d(r) < 0$: capital is abundant, lower $r$.
\item Bisection or Brent's method converges in $O(\log(1/\text{tol}))$ outer iterations.
\end{itemize}
\medskip
Each outer iteration requires solving Steps 1 and 2 from scratch --- the dominant cost.
\end{frame}

\begin{frame}{Implementation details}
\begin{itemize}
\item \textbf{Asset grid}: dense near the borrowing constraint where the policy bends; sparse for high $a$. Often log-spaced.
\medskip
\item \textbf{Initial value function}: zero, or the closed-form complete-markets value, or the value from the previous outer iteration.
\medskip
\item \textbf{Multigrid}: solve on a coarse grid first, then interpolate onto a finer grid as a warm start.
\medskip
\item \textbf{Acceleration}: Howard's policy iteration steps, modified policy iteration, endogenous grid method (Carroll, 2006).
\medskip
\item \textbf{Convergence criterion}: $\|v^{(n+1)} - v^{(n)}\|_\infty < \varepsilon_v$ inside; $|d(r)| < \varepsilon_r$ outside.
\end{itemize}
\end{frame}

\begin{frame}{Numerical challenges and the curse of dimensionality}
\begin{itemize}
\item For Aiyagari, the per-iteration state space is $\mathbb{R}_+ \times \mathcal{E}$ --- 2 continuous dimensions, manageable on a grid.
\medskip
\item Real-world quantitative HA models often have:
  \begin{itemize}
  \item Multiple assets (liquid + illiquid in HANK; portfolio choice).
  \item Multidimensional idiosyncratic states (income, unemployment, family structure).
  \item \emph{Aggregate} shocks $\to$ the cross-sectional distribution becomes part of the state (Krusell--Smith).
  \end{itemize}
\medskip
\item With $d$ continuous individual states and $N_a$ grid points per dimension, total points are $N_a^d$ --- the \textbf{curse of dimensionality}.
\medskip
\item This is where machine learning, and in particular deep neural networks, become indispensable.
\end{itemize}
\end{frame}

%=============================================================================
% PART III: WHY ML? MOTIVATING EXAMPLE ON OPTIMAL GROWTH
%=============================================================================

\section{Why ML? Motivating Example}

\begin{frame}{Why approximate the value function with a neural network?}
\begin{itemize}
\item Grid-based methods scale exponentially: $N_a^d$ points for $d$ continuous dimensions.
\medskip
\item Even modest models (income $\times$ wealth $\times$ housing $\times$ aggregate) become intractable.
\medskip
\item \textbf{Neural networks}:
  \begin{itemize}
  \item Universal approximators with smooth basis functions.
  \item Number of parameters scales (essentially) linearly in input dimension.
  \item Compatible with stochastic gradient training and GPU hardware.
  \end{itemize}
\medskip
\item Strategy: parameterise $V$ and $g$ by neural networks $\Gamma_{\gamma_v}$ and $\Gamma_{\gamma_g}$, then update $(\gamma_v,\gamma_g)$ by training.
\medskip
\item We illustrate first on a simple deterministic optimal growth model, then scale up.
\end{itemize}
\end{frame}

\begin{frame}{Optimal growth as a toy: NN approximations of $V$ and $g$}
The deterministic Bellman equation:
\[
V(k) \;=\; \max_{g}\!\left\{u(k,g(k)) + \beta\, V\!\big(f(k,g(k))\big)\right\}.
\]
\medskip
Parameterise:
\[
V(k) \approx \Gamma_{\gamma_v}(k), \qquad g(k) \approx \Gamma_{\gamma_g}(k).
\]
\medskip
\textbf{Algorithm sketch}:
\begin{itemize}
\item Sample states $k \sim d(k)$.
\item Compute target values $\hat V(k)$ from the right-hand side using current networks.
\item Update $\gamma_v$ to minimise $(\Gamma_{\gamma_v}(k) - \hat V(k))^2$.
\item Update $\gamma_g$ by gradient ascent on the per-state objective.
\item Iterate.
\end{itemize}
\end{frame}

\begin{frame}{One-step value update ($T_{\text{sim}}=1$)}
Given current policy $g^{(n-1)}$ and value $\Gamma_{\gamma_v}^{(n-1)}$, define the bootstrap target
\[
\hat V(k) \;=\; u\!\big(k,\, g^{(n-1)}(k)\big) + \beta\,\Gamma_{\gamma_v}^{(n-1)}\!\big(f(k, g^{(n-1)}(k))\big).
\]
\medskip
Then update by minimising the mean-squared Bellman residual:
\[
\gamma_v^{(n)} \;\in\; \arg\min_{\gamma_v}\; \mathbb{E}_{k\sim d}\Big[\big(\Gamma_{\gamma_v}(k) - \hat V(k)\big)^2\Big].
\]
\medskip
\begin{itemize}
\item Estimate the expectation by a Monte Carlo average over a batch $\{k_i\}_{i=1}^{B}$.
\item One \emph{epoch} = several gradient steps on the loss.
\item This is essentially \emph{fitted value iteration} (Bertsekas \& Tsitsiklis 1996), with a NN function class.
\end{itemize}
\end{frame}

\begin{frame}{One-step policy update}
Given the just-updated value $\Gamma_{\gamma_v}^{(n)}$, improve the policy by maximising the per-state RHS over the parameters of the policy network:
\[
\gamma_g^{(n)} \;\in\; \arg\max_{\gamma_g}\;
\mathbb{E}_{k\sim d}\Big[\,u\!\big(k,\Gamma_{\gamma_g}(k)\big) + \beta\,\Gamma_{\gamma_v}^{(n)}\!\big(f(k,\Gamma_{\gamma_g}(k))\big)\Big].
\]
\medskip
\begin{itemize}
\item Gradient with respect to $\gamma_g$ is computed automatically (backpropagation).
\item One step of Adam (or any first-order optimiser) on a Monte Carlo batch.
\item Alternating these updates is structurally identical to \textbf{actor--critic} reinforcement learning.
\end{itemize}
\end{frame}

\begin{frame}{Choosing the sampling distribution $d(k)$}
\textbf{Question:} on which states $k$ should we evaluate the loss?
\medskip
\begin{itemize}
\item \emph{Uniform} on the bounded grid: simple but wasteful (most weight on irrelevant states).
\medskip
\item \emph{Stationary distribution under the current policy}: focuses learning where the model actually spends time. Computed by simulation.
\medskip
\item \emph{Importance sampling}: re-weight to keep the gradient unbiased while sampling from a convenient proposal $q(k)$:
\[
\nabla_{\gamma}\mathbb{E}_{d}[L(k;\gamma)] = \mathbb{E}_q\!\left[\frac{d(k)}{q(k)}\nabla_\gamma L(k;\gamma)\right].
\]
\end{itemize}
\medskip
The choice of $d(k)$ has \emph{first-order} implications for solution accuracy: a network trained on the wrong support generalises poorly.
\end{frame}

\begin{frame}{Bias--variance tradeoff for $T_{\text{sim}}=1$}
\textbf{Pros of one-step bootstrapping}:
\begin{itemize}
\item Cheap: each evaluation costs a single network forward pass.
\item Low variance.
\end{itemize}
\medskip
\textbf{Cons}:
\begin{itemize}
\item High bias: the target $\hat V$ relies on the previous (imperfect) network.
\item Errors compound across iterations if the network is far from the true $V$.
\end{itemize}
\medskip
\textbf{Multi-step rollouts ($T_{\text{sim}}>1$)}:
\[
\hat V(k) \;=\; \sum_{t=0}^{T-1}\beta^t u(k_t, g(k_t)) + \beta^T \Gamma_{\gamma_v}\!(k_T).
\]
\begin{itemize}
\item Reduces bias (less reliance on the bootstrapped tail).
\item Increases variance and per-step cost.
\item Choose $T$ to trade off accuracy vs.\ wall time.
\end{itemize}
\end{frame}

\begin{frame}{Improvement: stationary-distribution sampling}
\begin{itemize}
\item Naive uniform sampling over a large bounded grid is wasteful: most points are never visited in equilibrium.
\medskip
\item Better: simulate the economy under the current policy until it converges to its (induced) stationary distribution, then sample $k$ from that distribution.
\medskip
\item Benefits:
  \begin{itemize}
  \item \emph{Sample efficiency}: training data concentrates where the model lives.
  \item \emph{Faster convergence}: gradients informative.
  \item \emph{LLN}: large simulation $\Rightarrow$ low-variance estimates of expectations.
  \end{itemize}
\medskip
\item Caveat: the sampling distribution \emph{depends on the current policy}, which itself is being trained --- creates a moving target. Standard fix: refresh the simulation every few epochs.
\end{itemize}
\end{frame}

\begin{frame}{An illustration: optimal growth under stochastic productivity}
\begin{center}
\includegraphics[width=0.55\textwidth]{optgrowth_sto}
\end{center}
\vspace{-0.5em}
\begin{itemize}
\item Stationary distribution of capital under the trained policy, vs.\ analytical benchmark.
\item NN-trained policy matches the benchmark closely after a few thousand epochs of alternating updates.
\end{itemize}
\end{frame}

\begin{frame}{What does RL/DL give us beyond grids?}
\begin{itemize}
\item \textbf{Function approximation that scales with data, not with grid size.} A 4-layer MLP with a few thousand parameters can represent value functions in $\mathbb{R}^{20}$ that no Cartesian grid could store.
\medskip
\item \textbf{Stochastic gradient training.} Each gradient step uses a small batch; we can train indefinitely on streaming Monte Carlo data.
\medskip
\item \textbf{Hardware acceleration.} Forward and backward passes are GPU-friendly matrix multiplications.
\medskip
\item \textbf{End-to-end differentiability.} Bellman residuals, Euler residuals, and even simulation steps can be differentiated through.
\medskip
\item These advantages compound when we move from optimal growth $\to$ Aiyagari $\to$ Krusell--Smith $\to$ HANK.
\end{itemize}
\end{frame}

%=============================================================================
% PART IV: EULER-EQUATION + ML ON OPTIMAL GROWTH
%=============================================================================

\section{Euler Equation + ML}

\begin{frame}{Euler equation as an alternative optimality condition}
For the optimal growth model with $u'(c) > 0$ and the budget constraint $c + k' = f(k)$, the consumer's first-order condition is the \textbf{Euler equation}:
\[
u'\!\big(c(k)\big) \;=\; \beta\, u'\!\big(c(k')\big)\, f'(k').
\]
\medskip
\begin{itemize}
\item Equivalent to the Bellman equation (under standard regularity).
\item Does not require representing $V$ explicitly: we can train only the policy network.
\item For the stochastic case with shock $z$:
\[
u'(c(k,z)) = \beta\,\mathbb{E}\big[u'(c(k',z'))\,(1+r(k',z'))\,\big|\,z\big].
\]
\end{itemize}
\end{frame}

\begin{frame}{From Euler residual to MSE loss}
Define the \textbf{Euler residual} at state $k$ as
\[
\varepsilon(k;\gamma_g) \;=\; 1 \,-\, \frac{\beta\,\mathbb{E}\!\left[u'(c'(k;\gamma_g))\,(1+r')\,\big|\,k\right]}{u'(c(k;\gamma_g))}.
\]
\medskip
The optimal policy makes $\varepsilon(k)\equiv 0$. Cast as supervised learning:
\[
\gamma_g^{*} \;\in\; \arg\min_{\gamma_g}\; \mathbb{E}_{k\sim d}\!\left[\varepsilon(k;\gamma_g)^2\right].
\]
\medskip
\begin{itemize}
\item Loss is \emph{differentiable} with respect to $\gamma_g$ via autograd.
\item Conditional expectation approximated by Monte Carlo over future shocks.
\item No explicit value function required.
\end{itemize}
\end{frame}

\begin{frame}{Minimising the Euler-residual loss}
\textbf{Algorithm}:
\begin{enumerate}
\item Initialise $\gamma_g$ randomly (e.g., Xavier/Glorot).
\item Sample a batch of states $\{k_i\}_{i=1}^B$ from $d(k)$.
\item For each $k_i$, compute $c_i = c(k_i;\gamma_g)$ and $k'_i = f(k_i) - c_i$.
\item Sample shocks $z'$ and compute the Monte Carlo expectation in $\varepsilon$.
\item Compute the squared residual loss and its gradient wrt $\gamma_g$.
\item Update $\gamma_g$ using Adam.
\item Repeat until convergence ($\sum_i \varepsilon_i^2 < $ tol).
\end{enumerate}
\medskip
\textbf{Why this is attractive}: \emph{first-order necessary conditions} are precisely what we want the policy to satisfy. The neural network is just a flexible function class.
\end{frame}

\begin{frame}[fragile]{Code 1/5: ValueNet architecture (PyTorch)}
\begin{lstlisting}
import torch
import torch.nn as nn

class ValueNet(nn.Module):
    """Two-hidden-layer MLP V_gamma(k)."""
    def __init__(self, n_hidden=64):
        super().__init__()
        self.net = nn.Sequential(
            nn.Linear(1, n_hidden),
            nn.ReLU(),
            nn.Linear(n_hidden, n_hidden),
            nn.ReLU(),
            nn.Linear(n_hidden, 1),
        )

    def forward(self, k):
        return self.net(k)
\end{lstlisting}
\medskip
A \emph{policy} network has identical structure but different weights $\gamma_g$.
\end{frame}

\begin{frame}[fragile]{Code 2/5: value update step (Bellman MSE)}
\begin{lstlisting}
def value_update_step(value_net, policy_net, k_batch,
                       beta, opt_v):
    """One gradient step on the Bellman MSE loss."""
    opt_v.zero_grad()
    c_now   = policy_net(k_batch)
    k_next  = production(k_batch) - c_now
    v_next  = value_net(k_next).detach()
    target  = utility(c_now) + beta * v_next
    v_now   = value_net(k_batch)
    loss    = ((v_now - target) ** 2).mean()
    loss.backward()
    opt_v.step()
    return loss.item()
\end{lstlisting}
\medskip
Note the \texttt{.detach()} on \texttt{v\_next}: we treat the bootstrap target as a constant.
\end{frame}

\begin{frame}[fragile]{Code 3/5: policy improvement (gradient ascent)}
\begin{lstlisting}
def policy_step(value_net, policy_net, k_batch,
                 beta, opt_g):
    """Maximise u(c) + beta V(k') with respect to policy params."""
    opt_g.zero_grad()
    c       = policy_net(k_batch)
    k_next  = production(k_batch) - c
    objective = (utility(c) + beta * value_net(k_next)).mean()
    loss = -objective                 # gradient *ascent*
    loss.backward()
    opt_g.step()
    return objective.item()
\end{lstlisting}
\medskip
The budget constraint $c+k'=f(k)$ is encoded inside the function so the policy net only outputs $c$.
\end{frame}

\begin{frame}[fragile]{Code 4/5: training loop}
\begin{lstlisting}
def train(value_net, policy_net, n_epochs, batch_size,
          beta, lr_v=1e-3, lr_g=1e-3):
    opt_v = torch.optim.Adam(value_net.parameters(), lr=lr_v)
    opt_g = torch.optim.Adam(policy_net.parameters(), lr=lr_g)

    history = {"v_loss": [], "obj": []}
    for epoch in range(n_epochs):
        k_batch = sample_states(batch_size)        # from d(k)
        v_loss  = value_update_step(value_net, policy_net,
                                     k_batch, beta, opt_v)
        obj     = policy_step(value_net, policy_net,
                               k_batch, beta, opt_g)
        history["v_loss"].append(v_loss)
        history["obj"].append(obj)
    return history
\end{lstlisting}
\end{frame}

\begin{frame}[fragile]{Code 5/5: visualisation}
\begin{lstlisting}
import matplotlib.pyplot as plt

def plot_results(value_net, policy_net):
    k = torch.linspace(0.1, 5.0, 100).reshape(-1, 1)
    with torch.no_grad():
        v = value_net(k).squeeze().numpy()
        c = policy_net(k).squeeze().numpy()

    fig, ax = plt.subplots(1, 2, figsize=(10, 4))
    ax[0].plot(k.numpy(), v); ax[0].set_title("V(k)")
    ax[1].plot(k.numpy(), c); ax[1].set_title("c(k)")
    plt.show()
\end{lstlisting}
\medskip
A few thousand epochs of alternating value/policy updates produce policies that match the analytical benchmark to within numerical tolerance.
\end{frame}

\begin{frame}{Actor--critic perspective: this is RL}
\begin{itemize}
\item The two-network alternating update is exactly the structure of \textbf{actor--critic} reinforcement learning:
\[
\text{Critic:}\quad \Gamma_{\gamma_v} \leftarrow \arg\min \;\big\|\Gamma_{\gamma_v} - \hat V\big\|^2,
\]
\[
\text{Actor:}\quad \Gamma_{\gamma_g} \leftarrow \arg\max \;\mathbb{E}\!\left[u + \beta V\right].
\]
\medskip
\item Difference vs.\ classical RL: in macro the environment (the model) is \emph{known}, so we can compute targets analytically rather than by Monte Carlo over an unknown $P(s'|s,a)$.
\medskip
\item This makes deep-learning macro \emph{model-based RL with a known model} --- often called ``deep dynamic programming''.
\end{itemize}
\end{frame}

%=============================================================================
% PART V: GLOBAL SOLUTION VIA DNN WITH FULL DISTRIBUTION DYNAMICS
%=============================================================================

\section{Global Solution via DNN}

\begin{frame}{Why (deep) RL helps}
Three reasons NN-based methods outperform grid methods on HA models:
\medskip
\begin{enumerate}
\item \textbf{Flexible function approximation.}\\
No structured Cartesian grid; networks handle high-dimensional inputs gracefully.
\medskip
\item \textbf{Sampling-based training.}\\
RL methods were designed for unknown environments; here we re-use the data-generation tools (rollouts, replay, importance sampling) even though our model is known. This avoids constructing exhaustive grids or closed-form distribution evolution.
\medskip
\item \textbf{Hardware acceleration.}\\
Forward/backward passes through neural networks are highly parallel matrix operations that GPUs and TPUs execute efficiently. Grid VFI is much harder to parallelise.
\end{enumerate}
\end{frame}

\begin{frame}{Global solution: full distribution as part of the state}
\begin{itemize}
\item In Aiyagari we exploited \emph{stationarity}: the distribution $\lambda$ was a fixed object and dropped out of the household state.
\medskip
\item Once we add aggregate shocks, $\lambda_t$ becomes part of the dynamic state. The household problem is
\[
V(a, z, \lambda) = \max_{c,a'}\!\left\{u(c) + \beta\,\mathbb{E}\!\left[V(a', z', \lambda')\,|\,z\right]\right\}.
\]
\item The law of motion for $\lambda$ is
\[
\lambda' \;=\; G(\lambda, Z, Z'),
\]
itself induced by the policy $a'(\cdot, \cdot, \lambda)$ and the shock kernel.
\item $\lambda$ is an infinite-dimensional object --- this is the heart of the curse of dimensionality in HA models with aggregate shocks.
\end{itemize}
\end{frame}

\begin{frame}{Household problem with an evolving distribution}
At the individual level, with state $(a,z,\lambda)$:
\[
V(a,z,\lambda) \;=\; \max_{a'}\!\left\{u\!\big((1+r(\lambda))a + w(\lambda)z - a'\big) + \beta\,\mathbb{E}\!\left[V(a',z',\lambda')\,|\,z\right]\right\}
\]
with $\lambda' = G(\lambda, Z, Z')$.
\medskip
\begin{itemize}
\item Two challenges:
  \begin{enumerate}
  \item How to represent $\lambda$ compactly enough that a NN can take it as input.
  \item How to evolve $\lambda$ forward without simulating $\sim 10^6$ agents.
  \end{enumerate}
\medskip
\item Three families of solutions:
  \begin{itemize}
  \item \emph{Moment approximations} (Krusell--Smith): summarise $\lambda$ by a few moments.
  \item \emph{Generalised moments} (DeepHAM): let the network learn the right summary statistics.
  \item \emph{Histogram transport} (DEQN): keep the full discretised distribution.
  \end{itemize}
\end{itemize}
\end{frame}

\begin{frame}{Approximating the distribution as a NN input}
\begin{center}
\includegraphics[width=0.55\textwidth]{app_dist}
\end{center}
\vspace{-0.5em}
\begin{itemize}
\item One option: discretise $\lambda$ on a fixed grid, feed the histogram heights to the NN.
\item Another: pass a finite vector of moments (means, variances, percentiles).
\item DeepHAM: \emph{learn} basis functions $Q_1,\ldots,Q_J$ and feed $\frac{1}{N}\sum_i Q_j(a_i)$ as moments.
\end{itemize}
\end{frame}

\begin{frame}{DNN representation of $V$ and the policy}
Parameterise
\[
V(S;\theta^V), \qquad c(S;\theta^c), \qquad S = (a, z, \lambda).
\]
\medskip
\begin{itemize}
\item Each network is a multilayer perceptron with parameters $\theta = (\text{weights}, \text{biases})$.
\item Inputs: individual state $(a,z)$ + summary of $\lambda$ (moments or generalised moments).
\item Outputs: scalar $V$, scalar consumption (or savings rate).
\medskip
\item Activations: ReLU, GELU, Mish --- all popular for HA-style problems.
\item Initialisation: Xavier/Glorot or Kaiming, depending on the activation.
\end{itemize}
\end{frame}

\begin{frame}{Update value function via simulation + supervised learning}
Given current policy $f_g^{(n)}$, generate $N_V$ initial states $\{S_{i,0}\}_{i=1}^{N_V}$ and simulate $T_V$ periods forward. Compute the empirical (truncated) discounted utility:
\[
\hat V(S_{i,0}) \;=\; \sum_{t=0}^{T_V-1} \beta^t \, u(c_{i,t}).
\]
\medskip
Train the value network by supervised regression:
\[
\theta^{V,(n+1)} \;\in\; \arg\min_{\theta^V}\;
\sum_{i=1}^{N_V} \big(V(S_{i,0};\theta^V) - \hat V(S_{i,0})\big)^2.
\]
\medskip
\begin{itemize}
\item This is sometimes called \emph{Monte Carlo value estimation}.
\item Larger $T_V$ reduces bias but increases variance and cost.
\end{itemize}
\end{frame}

\begin{frame}{Update policy with terminal value function}
Given the just-updated value network $V(\cdot;\theta^{V,(n+1)})$, choose policy parameters to maximise the truncated objective with terminal value:
\[
\theta^{c,(n+1)} \;\in\; \arg\max_{\theta^c}\;
\mathbb{E}\!\left[\sum_{t=0}^{T_g-1}\beta^t u(c_{i,t}) + \beta^{T_g}\, V(S_{i,T_g};\theta^{V,(n+1)})\right]
\]
subject to the budget constraint
\[
c_{i,t} = (1+r(\lambda_t))\,a_{i,t} + w(\lambda_t)\,z_{i,t} - a'_{i,t},
\quad a'_{i,t} = c(S_{i,t};\theta^c).
\]
\medskip
This is a long-horizon \emph{differentiable rollout}: gradients propagate through the simulated path.
\end{frame}

\begin{frame}{Transition function $G$ as a matrix}
On a discretised asset grid $\{a_1,\ldots,a_{N_a}\}$, the policy $a'(\cdot)$ together with the shock kernel induces a Markov matrix
\[
\hat G_{i,j} \;=\; \Pr\!\big(a' \in \phi_j^a \,\big|\, a \in \phi_i^a\big),
\]
where $\phi_j^a$ are the discretisation cells.
\medskip
\begin{itemize}
\item $\hat G$ has $O(N_a^2)$ entries; sparse if the policy moves between adjacent cells.
\item Updating $\lambda$ one period: $\lambda^{(t+1)} = \hat G^\top \lambda^{(t)}$.
\item The same matrix is used in classical Aiyagari computation; here we re-use it inside the differentiable training loop.
\end{itemize}
\end{frame}

\begin{frame}{Conditional and unconditional transition kernels}
\textbf{Unconditional:}
\[
\lambda'(z' \in Z',\, a' \in A')
\;=\;
\int_{Z\times A} \mathbb{1}\!\big\{a'(a,z,\lambda)\in A',\; z'(a,z,\lambda)\in Z'\big\}\, d\lambda(z,a).
\]
\medskip
\textbf{Conditional (given current $(a,z)$):}
\[
\lambda'(z' \in Z',\, a' \in A' \mid z\in Z, a\in A)
=
\int_{Z\times A}\!\!\mathbb{1}\!\{\cdots\}\, d\lambda(z,a) \;/\; \lambda(Z\times A).
\]
\medskip
\begin{itemize}
\item Both kernels live in the Wasserstein space of probability measures on $S$.
\item The induced \emph{stochastic process on distributions} is the analytical object that DeepHAM and DEQN approximate by neural networks.
\end{itemize}
\end{frame}

\begin{frame}{Discretised transition: i.i.d.\ shocks, monotone policy}
When $z$ is i.i.d.\ and the policy $a'(a,z)$ is monotone in $a$ for each $z$, we can solve for transition cell boundaries explicitly:
\[
\hat G^{\lambda,\eta_x}_{i,j}
\;=\; \Pr\!\big(\phi_j^a \ni a_+ \leq \bar\phi_j^a \,\big|\, a \in \phi_i^a\big)
\;=\; F_z(\bar z_{i,j}) - F_z(\underline z_{i,j}),
\]
where $(\underline z_{i,j}, \bar z_{i,j})$ are the income shock thresholds that map $\phi_i^a$ to the boundaries of $\phi_j^a$ under the policy.
\medskip
\begin{itemize}
\item Closed-form bin probabilities $\Rightarrow$ no Monte Carlo noise in the $\lambda$-update.
\item Differentiability is delicate: the indicator functions inside the kernel are non-smooth, and need to be replaced by smooth approximations (e.g.\ sigmoids with sharp slopes) for gradient propagation.
\end{itemize}
\end{frame}

%=============================================================================
% PART VI: KRUSELL--SMITH (1998)
%=============================================================================

\section{Krusell--Smith (1998)}

\begin{frame}{Krusell--Smith (1998): aggregate uncertainty}
\begin{itemize}
\item Aiyagari ignores aggregate fluctuations. Krusell--Smith (\emph{JPE} 1998) adds an \textbf{aggregate TFP shock} $Z_t$ to the production technology.
\medskip
\item Setup:
  \begin{itemize}
  \item Continuum of households with idiosyncratic employment $\epsilon \in \{0,1\}$.
  \item Aggregate state $Z_t \in \{Z_g, Z_b\}$, two-state Markov.
  \item Joint $(\epsilon, Z)$ Markov chain (correlated unemployment and recessions).
  \item Inelastic labour, single risk-free bond, borrowing constraint.
  \end{itemize}
\medskip
\item The cross-sectional distribution $\Gamma_t$ over $(a,\epsilon)$ now \textbf{moves stochastically over time}.
\end{itemize}
\end{frame}

\begin{frame}{KS distributions, firms, prices}
\begin{itemize}
\item Aggregate state: $(Z, \Gamma)$ where $\Gamma$ is the joint distribution over individual asset and employment states.
\medskip
\item Firm: $Y_t = Z_t F(K_t, \bar L)$ with $K_t = \int a\,d\Gamma_t$ and $\bar L$ aggregate labour.
\medskip
\item Factor prices:
\[
R_t \;=\; Z_t F_K(K_t, \bar L) - \delta, \qquad
w_t \;=\; Z_t F_L(K_t, \bar L).
\]
\medskip
\item Both $R_t$ and $w_t$ depend on the aggregate state through $(Z_t, K_t)$, which depends on the entire \emph{history} of the economy.
\end{itemize}
\end{frame}

\begin{frame}{Household problem with the full distribution as state}
The HH Bellman equation is now
\[
V(a_i, y_i, Z, \Gamma) \;=\; \max_{c_i, a'_i}
\big\{u(c_i) + \beta\,\mathbb{E}\!\left[V(a'_i, y'_i, Z', \Gamma')\,\big|\,y_i, Z\right]\big\}
\]
subject to
\[
a'_{i} = w(Z,K)\,\bar\ell\, y_i + R(Z,K)\,a_i - c_i, \qquad K = \int a\,d\Gamma,
\]
and the law of motion $\Gamma' = G(\Gamma, Z, Z')$.
\medskip
\begin{itemize}
\item \textbf{Curse of dimensionality}: $\Gamma$ is infinite-dimensional.
\item Without simplification, the household must \emph{forecast the entire distribution} to make optimal decisions.
\end{itemize}
\end{frame}

\begin{frame}{State reduction via finite moments}
Krusell and Smith (1998) propose: replace $\Gamma$ with a finite vector of moments
\[
m \;=\; (m_1, m_2, \ldots, m_J), \qquad m_j \;=\; \int a^j\, d\Gamma.
\]
\medskip
\begin{itemize}
\item Most often: $J = 1$ --- only the first moment $m_1 = K = \int a\,d\Gamma$.
\medskip
\item This dramatically reduces the household state from $(a,y,Z,\Gamma)$ to $(a,y,Z,K)$.
\medskip
\item Approximation accuracy is excellent for many calibrations: agents care primarily about \emph{average} capital because that determines $r$ and $w$.
\medskip
\item ``Approximate aggregation'' (KS): higher moments rarely matter for the policy.
\end{itemize}
\end{frame}

\begin{frame}{Postulated log-linear law of motion}
Households need to forecast $K_{t+1}$ given $(Z_t, K_t)$. Postulate
\[
\log K_{t+1} \;=\; A(Z_t) + B(Z_t)\,\log K_t,
\]
with state-dependent intercept $A(Z)$ and slope $B(Z)$.
\medskip
\begin{itemize}
\item Two unknowns per aggregate state: $\{A(Z_g), B(Z_g), A(Z_b), B(Z_b)\}$.
\item Households use this rule to compute $\mathbb{E}[V(a',y',Z',K')\,|\,y,Z]$.
\item In equilibrium, the postulated rule must match the rule \emph{generated} by the household policy under simulation.
\end{itemize}
\end{frame}

\begin{frame}{KS algorithm: steps 1--2}
\textbf{Step 1.} Specify the functional form for the law of motion. Initialise $\{A(Z), B(Z)\}$.
\medskip

\textbf{Step 2.} Given the guessed law of motion, solve the household DP problem with state $(a,y,Z,K)$:
\[
V(a,y,Z,K) = \max_{a'}\!\big\{u(c) + \beta\,\mathbb{E}\!\left[V(a',y',Z',K'(Z))\,|\,y,Z\right]\big\}
\]
with $K'(Z) = \exp\!\big(A(Z) + B(Z)\log K\big)$.
\medskip
\begin{itemize}
\item This is a 4-dimensional VFI problem (continuous in $a$ and $K$, discrete in $y$ and $Z$).
\item Standard tools: tensor-product grids, multivariate interpolation, Howard improvement.\footnote{\scriptsize Classical HA computation methods are covered in detail in the video series at \url{https://space.bilibili.com/2142649036/lists/7180709}.}
\end{itemize}
\end{frame}

\begin{frame}{KS algorithm: steps 3--5}
\textbf{Step 3.} Simulate the economy under the household policy: $T \approx 10{,}000$ periods, $N \approx 10{,}000$ agents. Record the time series $\{K_t, Z_t\}$.
\medskip

\textbf{Step 4.} Run state-dependent OLS on the simulated $\{K_t, Z_t\}$:
\[
\log K_{t+1} = \hat A(Z_t) + \hat B(Z_t)\,\log K_t + e_t.
\]
\medskip

\textbf{Step 5.} Update $\{A(Z), B(Z)\} \leftarrow \{\hat A(Z), \hat B(Z)\}$ and return to Step 2. Iterate until $|A_{\text{new}} - A_{\text{old}}| < \varepsilon$.
\medskip
\textbf{Convergence criterion}: in-sample $R^2$ of the regression close to $1.0$ (often $> 0.9999$ in KS!).
\end{frame}

\begin{frame}{Limitations of moment approximation}
\begin{itemize}
\item Although remarkably accurate \emph{on average}, the moment approximation has weaknesses:
\medskip
\begin{itemize}
\item \emph{Specific functional form} for the law of motion is hard to choose ex ante.
\item \emph{Higher-order properties} of the distribution (welfare, inequality measures) require more moments.
\item \emph{Optimal-policy analysis} (designing fiscal/monetary intervention) needs derivatives of the value function with respect to the full distribution, not just moments.
\end{itemize}
\medskip
\item \textbf{The natural alternative}: let a neural network \emph{learn} which features of $\Gamma$ matter.
\medskip
\item This is exactly the DeepHAM / DEQN approach.
\end{itemize}
\end{frame}

%=============================================================================
% PART VII: DEEPHAM
%=============================================================================

\section{DeepHAM}

\begin{frame}{DeepHAM (Han, Yang \& E, 2025)}
\begin{itemize}
\item DeepHAM solves the (finite-$N$, large-$N$) Krusell--Smith problem as a \textbf{Markov Nash equilibrium} among $N$ agents.
\medskip
\item Two neural networks:
  \begin{itemize}
  \item Value function $V(s;\theta^V)$ for individual states.
  \item Policy $\mathcal{C}(s;\theta^C)$ for individual savings/consumption.
  \end{itemize}
\medskip
\item Training alternates:
  \begin{enumerate}
  \item \textbf{Value step}: simulate under current policy, compute Monte Carlo target, do supervised regression on $\theta^V$.
  \item \textbf{Policy step}: solve agent $i$'s problem with all other agents fixed at the previous policy --- \emph{fictitious play}.
  \end{enumerate}
\medskip
\item This decomposition makes the high-dimensional joint problem tractable.
\end{itemize}
\end{frame}

\begin{frame}{DeepHAM: algorithm outline}
\textbf{Initialisation ($k=0$).} Initial policy $\mathcal{C}^{(0)}$, networks $\theta^{V,(0)}, \theta^{C,(0)}$.
\medskip

\textbf{$k$-th iteration:}
\begin{enumerate}
\item Prepare the stationary distribution $\mu(\mathcal{C}^{(k-1)})$ under the previous policy.
\medskip
\item Update the value network by supervised learning:
\[
\theta^{V,(k)} \in \arg\min_{\theta^V}\;
\mathbb{E}_\mu\!\left[\big(V(s_i;\theta^V) - \hat V_i\big)^2\right],
\]
with target $\hat V_i = \sum_{\tau=0}^{T_v-1}\beta^\tau u(c_\tau^i)$.
\medskip
\item Update the policy network by fictitious play:
\[
\theta^{C,(k)} \in \arg\max_{\theta^C}\;
\mathbb{E}_\mu\!\left[\sum_{t=0}^{T-1}\beta^t u(c_{i,t}) + \beta^T V(s_{i,T};\theta^{V,(k)})\right].
\]
\end{enumerate}
\end{frame}

\begin{frame}{DeepHAM: value function learning}
\begin{itemize}
\item Simulate the economy for $T_v$ periods under the current policy $\mathcal{C}^{(k-1)}$.
\medskip
\item For each sampled initial state $s_i$, record the realised consumption path and compute the discounted utility target
\[
\hat V_i \;=\; \sum_{\tau=0}^{T_v-1} \beta^{\tau}\, u(c_{i,\tau}).
\]
\medskip
\item Train $V(\cdot;\theta^V)$ as a regression on $\{(s_i, \hat V_i)\}$.
\medskip
\item This is exactly fitted-value Monte Carlo: standard supervised deep learning, no special algorithms required.
\end{itemize}
\end{frame}

\begin{frame}{DeepHAM: policy iteration via fictitious play}
\begin{itemize}
\item In round $\ell+1$, all agents' policies $\mathcal{C}^{(k,\ell)}$ are known. Agent $i$ improves \emph{best-responding} to the previous policy of all others:
\[
\mathcal{C}^{(k,\ell+1)} \in \arg\max_{\mathcal{C}}\;
\mathbb{E}_\mu\!\left[\,\tilde U_{i,T} + \beta^T V^{(k)}(s_{i,T})\,\right].
\]
\medskip
\item Optimisation is over the parameters of the policy network and runs by gradient ascent on a Monte Carlo sample of paths.
\medskip
\item \textbf{Why fictitious play?}
  \begin{itemize}
  \item Decouples the high-dimensional fixed-point problem.
  \item Empirically robust: even in non-convex problems, alternating best response converges to a Nash equilibrium under regularity conditions.
  \end{itemize}
\end{itemize}
\end{frame}

\begin{frame}{DeepHAM: computational graph}
\begin{center}
\includegraphics[width=0.85\textwidth]{KSwithDL_architecture_price_new}
\end{center}
\vspace{-0.4em}
\begin{itemize}
\item Inputs: individual state $s = (a_i, y_i, Z, \Gamma)$, with $\Gamma$ summarised by generalised moments.
\item Layers: feedforward MLP with multiple hidden layers.
\item Outputs: scalar consumption decision; budget constraint enforces $a' = R a + w y - c$.
\end{itemize}
\end{frame}

\begin{frame}{DL-based generalised moments}
\begin{itemize}
\item Instead of a fixed list of moments $\{m_1, m_2, \ldots\}$, DeepHAM lets the network \emph{learn} symmetry-preserving sufficient statistics:
\[
\bar Q_j(\Gamma) \;=\; \frac{1}{N} \sum_{i=1}^{N} Q_j(a_i),
\]
where $Q_1, \ldots, Q_J$ are basis functions parameterised by additional NN layers.
\medskip
\item Special case $Q_j(a) = a^j$ recovers the classical moments.
\medskip
\item Trained jointly with $V$ and $\mathcal C$, the learned moments adapt to the economic problem at hand.
\medskip
\item Symmetry across agents (the population is exchangeable) is enforced by averaging $Q_j$ across the cross section.
\end{itemize}
\end{frame}

\begin{frame}{Interpretation of generalised moments}
\begin{center}
\includegraphics[width=0.85\textwidth]{KS_0fm1gm_v1_valuegmbasis_sub}
\end{center}
\vspace{-0.4em}
\begin{itemize}
\item Each panel: a learned basis function $Q_j(a)$ as a function of individual asset $a$.
\item The first basis function looks roughly linear ($\sim$ first moment); subsequent functions capture concavities and tail behaviour.
\item The network discovers, automatically, the features of $\Gamma$ that are relevant for forecasting.
\end{itemize}
\end{frame}

\begin{frame}{DeepHAM output figures}
\begin{center}
\includegraphics[width=0.42\textwidth]{deepham}\hfill
\includegraphics[width=0.42\textwidth]{ddp_training_metrics}\\[0.5em]
\includegraphics[width=0.42\textwidth]{scatter_policy_c}\hfill
\includegraphics[width=0.42\textwidth]{scatter_policy_a1}
\end{center}
\vspace{-0.5em}
\begin{itemize}
\item \emph{Top left}: simulated cross section of capital under the trained policy.
\item \emph{Top right}: training loss curves.
\item \emph{Bottom}: consumption and savings policies as functions of individual state, scattered across cross-section samples.
\end{itemize}
\end{frame}

%=============================================================================
% PART VIII: DEEP EQUILIBRIUM NETS (DEQN)
%=============================================================================

\section{Deep Equilibrium Nets (DEQN)}

\begin{frame}{Deep Equilibrium Nets (Azinovic, Gaegauf, Scheidegger 2022)}
\begin{itemize}
\item DEQN solves the same KS economy with a \textbf{single} neural network, no value function.
\medskip
\item The single network represents the policy:
\[
\sigma_\theta(a, e, Z, K) \in (0,1)
\;\Rightarrow\;
a' = \underline a + \sigma_\theta\cdot(\text{cash on hand} - \underline a).
\]
\medskip
\item The training objective is the \textbf{Euler equation residual}, not a Bellman residual.
\medskip
\item No value network, no fictitious play; just supervised learning on the Euler condition.
\medskip
\item Key innovation: a \emph{cloud} of parallel economies provides a rich training distribution.
\end{itemize}
\end{frame}

\begin{frame}{DEQN: policy network architecture}
\begin{itemize}
\item Single multilayer perceptron $\sigma_\theta : \mathbb{R}^4 \to (0,1)$.
\medskip
\item Inputs: $(a, e, Z, K)$ where $K = \int a\, d\mu$ is aggregate capital.
\medskip
\item Architecture used in the original paper:
  \begin{itemize}
  \item 3 hidden layers, 64 units each.
  \item Mish activation function.
  \item Sigmoid output for the savings rate.
  \end{itemize}
\medskip
\item Sigmoid output enforces the budget constraint at the network level: the implied $a'$ always lies within $[\underline a, \text{cash on hand}]$.
\medskip
\item Contrast with DeepHAM: DeepHAM uses both $V_{\text{NN}}$ and $\mathcal C_{\text{NN}}$; DEQN uses only the policy network.
\end{itemize}
\end{frame}

\begin{frame}{DEQN: Euler equation as the loss function}
The household optimality condition is the Euler equation:
\[
u'(c_t) \;=\; \beta\,\mathbb{E}_t\!\left[(1+r_{t+1})\,u'(c_{t+1})\right].
\]
\medskip
The \textbf{Euler residual}:
\[
\varepsilon \;=\; 1 \,-\, \frac{\beta\,\mathbb{E}\!\left[(1+r')\,u'(c')\right]}{u'(c)}.
\]
\medskip
The DEQN loss aggregates squared residuals across a sample of agents:
\[
\mathcal{L}(\theta) \;=\; \frac{1}{M}\sum_{m=1}^{M} \varepsilon_m^2.
\]
\medskip
\textbf{Why Euler equations?} They are necessary and sufficient for optimality (under regularity); driving $\varepsilon\to 0$ trains the optimal policy without ever computing $V$.
\end{frame}

\begin{frame}{DEQN: cloud simulation}
\textbf{Problem.} To evaluate the Euler residual we need $(K, Z)$, but $K$ depends on the cross-sectional distribution $\mu$, which itself evolves endogenously under the trained policy.
\medskip

\textbf{Solution.} Simulate $N$ \emph{independent copies} of the entire economy in parallel:
\medskip
\begin{itemize}
\item Epoch 0: all $N$ economies start at the steady-state distribution $\mu_{ss}$ from a sequence-Jacobian (SSJ) calibration.
\item As training progresses, each economy evolves under its own TFP path $\{Z_n\}$, so the set $\{K_n\}$ spreads out across a realistic range.
\item Result: an economically consistent and \emph{diverse} batch on which the Euler loss is meaningful.
\end{itemize}
\medskip
Compared to DeepHAM (one economy, $N$ agents, $T_v$-step rollout), DEQN runs $N$ economies, each with its own $\mu_n$, evaluating the Euler equation directly each step.
\end{frame}

\begin{frame}{DEQN: distribution transport}
At each epoch, $\mu_{n,t}$ is a histogram over $(a_j, e_k)$. Steps to advance one period:
\medskip

\begin{enumerate}
\item Compute savings $a'_{j,k}$ for every grid point under the current policy.
\item \textbf{Linear interpolation transport}: split mass $\mu(a_j, e_k)$ between the two adjacent grid points with weights $w_{\ell+1}, w_\ell$.
\item \textbf{Markov transition for employment}: distribute mass to $(a_\ell, e')$ weighted by the transition matrix $\Pi_{e\to e'}$.
\item Update $K_{n,t+1} = \sum_{j,k} a_j \cdot \mu_{n,t+1}(a_j, e_k)$.
\end{enumerate}
\medskip

The transport step runs inside \texttt{torch.no\_grad()} so that gradients only propagate through the Euler residual loss, not through the simulation step itself.
\end{frame}

\begin{frame}{DEQN: algorithm outline}
\textbf{Initialisation.} Solve the steady state via SSJ to get $K_{ss}$, asset/employment grids, $\mu_{ss}$, and the Markov transition $\Pi$.
\medskip

\textbf{For each training epoch $t \in \{1,\ldots,T\}$:}
\begin{enumerate}
\item Compute $K_n = \int a\,d\mu_n$ and read $Z_n$ for each cloud copy $n$.
\item Sample $M$ random $(a, e)$ per copy.
\item Forward pass: evaluate $\sigma_\theta \to (a', c)$.
\item Compute Euler residuals $\varepsilon^2$ via expectations over $(Z', e')$.
\item Backward pass: Adam update on $\theta$.
\item Transport: $\mu_n \to \mu_n'$ inside \texttt{no\_grad()}; draw next $Z_n$.
\end{enumerate}
\medskip

\textbf{Convergence.} Loss decreases over $\sim 2000$ epochs; Euler errors below $10^{-3}$ indicate the policy solves the model to high accuracy.
\end{frame}

\begin{frame}{DEQN vs.\ DeepHAM: comparison}
\begin{center}
\small
\begin{tabular}{lll}
\toprule
 & \textbf{DeepHAM} & \textbf{DEQN}\\
\midrule
Networks & Value $+$ policy (2) & Policy only (1)\\
Optimality & Bellman / fictitious play & Euler residual\\
Distribution rep.\ & Generalised moments & Full histogram\\
Training data & $N$ agents, 1 economy & $N$ economies, cloud sim\\
Forward sim & Yes, $T_v$ steps & Yes, 1 step / epoch\\
Distribution evolution & Implicit via moments & Linear-interp transport\\
Scalability & Moments fixed; NN width grows & Histogram sized to grid\\
Differentiability & Through truncated rollout & Euler loss only\\
\bottomrule
\end{tabular}
\end{center}
\medskip
\begin{itemize}
\item Both avoid the curse of dimensionality of the underlying KS problem.
\item DeepHAM has richer distributional info via learned moments; DEQN has a simpler architecture and direct Euler-based training.
\end{itemize}
\end{frame}

\begin{frame}{Recap: the modern HA toolbox}
\begin{itemize}
\item Aiyagari (1994): a non-trivial nested fixed point, classically solved by VFI + bisection on $r$.
\medskip
\item Krusell--Smith (1998): adds aggregate uncertainty, classically tamed by bounded-rationality moments.
\medskip
\item DeepHAM: replaces moments with \emph{learned} sufficient statistics; alternates supervised value learning and fictitious-play policy improvement.
\medskip
\item DEQN: drops the value function entirely; trains a single policy network on Euler residuals over a cloud of simulated economies.
\medskip
\item Both deep-learning methods generalise straightforwardly to:
  \begin{itemize}
  \item Multiple assets (HANK)
  \item OLG with aggregate shocks (next part)
  \item Mean-Field Games in continuous time (Part X)
  \end{itemize}
\end{itemize}
\end{frame}

%=============================================================================
% PART IX: OLG WITH AGGREGATE UNCERTAINTY
%=============================================================================

\section{OLG with Aggregate Uncertainty}

\begin{frame}{From infinite-horizon to overlapping generations}
\begin{itemize}
\item In Aiyagari/KS, agents are infinitely lived. Many quantitative questions require finite life cycles:
  \begin{itemize}
  \item Social security and pension reform
  \item Retirement savings and 401(k)
  \item Lifecycle wealth accumulation
  \item Education financing and student debt
  \end{itemize}
\medskip
\item \textbf{OLG model} (Diamond, Auerbach--Kotlikoff): agents live $J$ periods (or cohorts), are born, age, and die.
\medskip
\item State vector now includes the \emph{age} of each cohort, plus its asset/income distribution.
\medskip
\item Standard OLG with aggregate shocks remains computationally hard --- the per-cohort distribution is itself an infinite-dimensional state.
\end{itemize}
\end{frame}

\begin{frame}{OLG with aggregate uncertainty: state vector}
With $J$ cohorts indexed by age $j = 1, \ldots, J$:
\medskip
\begin{itemize}
\item Cohort $j$ has its own asset distribution $\Gamma_j$ over $(a_j, y_j)$.
\medskip
\item Aggregate state: $S = (Z, \Gamma_1, \Gamma_2, \ldots, \Gamma_J)$.
\medskip
\item Aggregate capital and labour:
\[
K_t = \sum_{j=1}^{J} \int a\,d\Gamma_{j,t}, \qquad
L_t = \sum_{j=1}^{J} \int \bar\ell_j\, y\, d\Gamma_{j,t},
\]
where $\bar\ell_j$ is the cohort-specific labour endowment (often hump-shaped).
\medskip
\item Lifecycle Bellman: $V_j(a_j, y_j, S) = \max\,\{u(c) + \beta\, \mathbb{E}[V_{j+1}(a',y',S')]\}$, with $V_{J+1}\equiv 0$ (or a terminal bequest function).
\end{itemize}
\end{frame}

\begin{frame}{Why DL helps for OLG with aggregate shocks}
\begin{itemize}
\item Without aggregate shocks: standard backward induction over cohorts works.
\medskip
\item \textbf{With} aggregate shocks: each cohort's continuation value depends on the future joint distribution of all cohorts $\Rightarrow$ exponential state-space blow-up.
\medskip
\item Deep neural networks let us:
  \begin{itemize}
  \item Approximate $V_j(a, y, S)$ with a single MLP per cohort (or one MLP with cohort as an input).
  \item Encode $S$ via learned summary statistics (analogous to DeepHAM moments).
  \item Train across cohorts in parallel on a GPU.
  \end{itemize}
\medskip
\item Recent reference: Brumm, Kubler \& Scheidegger (2023), ``Economics-inspired neural networks with stabilizing homotopies''.
\end{itemize}
\end{frame}

\begin{frame}{Homotopy-stabilised training (Brumm--Kubler--Scheidegger 2023)}
\begin{itemize}
\item OLG with aggregate uncertainty is hard to train: borrowing constraints, occasionally-binding bond limits, and discontinuities in the policy create unstable gradients.
\medskip
\item \textbf{Stabilising homotopy}: introduce constraints \emph{gradually} during training.
  \begin{itemize}
  \item Start with a relaxed (smooth) version of the model: no binding constraint, smooth penalties.
  \item Slowly tighten the constraints toward their true value as training progresses.
  \item By the time the constraints are at their true level, the network is already in a good basin.
  \end{itemize}
\medskip
\item Empirically, homotopy is essential for getting OLG-with-aggregate-shocks training to converge.
\medskip
\item Implementation often combines JAX (for automatic differentiation through long simulation rollouts) with the Adam optimiser and a multistep homotopy schedule.
\end{itemize}
\end{frame}

\begin{frame}{Lab12\_OLG.ipynb preview}
\begin{itemize}
\item \texttt{Notebooks/Lab12\_OLG.ipynb} is a JAX-based implementation of the Brumm--Kubler--Scheidegger OLG model with stabilising homotopies.
\medskip
\item Highlights:
  \begin{itemize}
  \item 7 cohorts spanning a 72-period life cycle, life-cycle labour endowment.
  \item Aggregate TFP shocks following an AR(1).
  \item Capital adjustment costs.
  \item State-dependent depreciation (TFP-tied).
  \item Two-stage homotopy: capital-only model first, then add bond.
  \end{itemize}
\medskip
\item Use this notebook as an \textbf{optional follow-up} after Lab06: students who finish the KS lab early can explore OLG.
\medskip
\item In production research, the same toolkit handles multi-asset HANK and asset-pricing models.
\end{itemize}
\end{frame}

%=============================================================================
% PART X: MEAN-FIELD GAMES
%=============================================================================

\section{Mean-Field Games}

\begin{frame}{Mean-Field Games: motivation}
\begin{itemize}
\item All the models we've seen so far are \emph{finite or large-but-finite} populations of agents (Aiyagari has a continuum but discretises numerically).
\medskip
\item \textbf{Mean-Field Games} (MFG) provide a continuous-time, continuum-of-agents formulation.
\medskip
\item Why care?
  \begin{itemize}
  \item Continuous time often \emph{simplifies} the math: PDEs replace functional equations.
  \item Powerful new toolbox for HA models with idiosyncratic and aggregate shocks.
  \item Tight link to optimal transport, stochastic analysis, and modern PDE theory.
  \item Recent breakthroughs in \emph{deep learning for MFG} (Carmona--Lauri\`ere, Achdou--Lasry--Lions, Han et al.).
  \end{itemize}
\medskip
\item This part is an introduction; you can think of MFG as the continuum limit of Aiyagari/KS.
\end{itemize}
\end{frame}

\begin{frame}{The core idea: continuum + mean-field feedback}
\begin{itemize}
\item Imagine a \emph{continuum} of identical agents, each solving an individual control problem.
\medskip
\item Each agent's payoff depends on the \emph{distribution} of all other agents (the ``mean field'').
\medskip
\item The distribution is itself determined by the optimal behaviour of the agents.
\medskip
\item This circular dependence is the defining feature of MFG --- and is exactly the same fixed point as in HA models, but stated in continuous time.
\medskip
\item Two mathematical objects describe the equilibrium:
  \begin{itemize}
  \item A backward Hamilton--Jacobi--Bellman (HJB) equation for the individual value function.
  \item A forward Kolmogorov forward / Fokker--Planck (KFE) equation for the distribution.
  \end{itemize}
\end{itemize}
\end{frame}

\begin{frame}{A brief history}
\begin{itemize}
\item \textbf{Engineering origins} (mid-2000s):\\
  Caines, Huang \& Malham\'e developed ``Nash certainty equivalence'' for large-population games.
\medskip
\item \textbf{Mathematical formalisation}:\\
  Lasry \& Lions (2006--2007) coined the term ``Mean-Field Games'' and developed the rigorous PDE theory.
\medskip
\item \textbf{Economic adoption}:\\
  Achdou--Han--Lasry--Lions--Moll (2014, 2022 \emph{REStud}) translated the framework into HA macroeconomics --- continuous-time Aiyagari--Bewley--Huggett.
\medskip
\item \textbf{Deep-learning era} (2020s):\\
  Carmona--Lauri\`ere, Han--Hu, Achdou--Pequito, $\ldots$ use deep neural networks to solve MFGs in high dimensions.
\end{itemize}
\end{frame}

\begin{frame}{The two-equation MFG system}
\textbf{Hamilton--Jacobi--Bellman (backward, individual)}:
\[
-\partial_t v(x,t) + \rho v(x,t) \;=\; H\!\big(x,\nabla v(x,t),m(\cdot,t)\big),
\]
where $H$ is the Hamiltonian of the individual control problem and $m(\cdot,t)$ is the population distribution at time $t$.
\medskip

\textbf{Kolmogorov forward / Fokker--Planck (forward, aggregate)}:
\[
\partial_t m(x,t) \;=\; -\nabla\!\cdot\!\big(s^*(x,t,m)\, m(x,t)\big) + \tfrac{1}{2}\Delta\!\big(\sigma^2(x,t)\,m(x,t)\big),
\]
where $s^*$ is the optimal drift derived from the HJB.
\medskip
\begin{itemize}
\item HJB and KFE are \emph{coupled}: the optimal control depends on $m$, and $m$ evolves under the optimal control.
\item In equilibrium, both equations must hold simultaneously.
\end{itemize}
\end{frame}

\begin{frame}{Why HJB is familiar, why KFE is new}
\begin{itemize}
\item \textbf{HJB} is the continuous-time limit of the Bellman equation. Macroeconomists have seen it many times: Merton's portfolio choice, deterministic optimal growth, $\ldots$
\medskip
\item \textbf{KFE / Fokker--Planck} describes the time evolution of a probability density under stochastic dynamics.
  \begin{itemize}
  \item Originated in physics: Fokker (1914), Planck (1917) --- diffusion of particles.
  \item Formalised by Kolmogorov (1931).
  \item Standard in physics and engineering, only recently mainstream in macro.
  \end{itemize}
\medskip
\item Together they describe the \emph{joint dynamics of optimal individual behaviour and the cross-sectional distribution} --- exactly what HA models need.
\end{itemize}
\end{frame}

\begin{frame}{HJB review: deterministic optimal control}
A generic optimal control problem:
\[
v(x_0) \;=\; \max_{\alpha(\cdot)} \int_0^\infty e^{-\rho t}\, r(x_t, \alpha_t)\, dt
\quad\text{s.t.}\quad \dot x_t = f(x_t, \alpha_t),\; x_0 \text{ given}.
\]
\medskip

The corresponding HJB equation:
\[
\rho\, v(x) \;=\; \max_{\alpha}\left\{r(x,\alpha) + v'(x)\,f(x,\alpha)\right\}.
\]
\medskip

\textbf{Interpretation}:
\begin{itemize}
\item LHS: ``required return'' on the value function.
\item RHS: instantaneous payoff plus the marginal value of the state moving in the chosen direction.
\item First-order condition over $\alpha$: $r_\alpha + v'(x) f_\alpha = 0$.
\end{itemize}
\end{frame}

\begin{frame}{From discrete-time Bellman to continuous-time HJB}
Start with the discrete-time Bellman over a small interval $\Delta$:
\[
v(x) \;\approx\; \max_{\alpha}\!\left\{r(x,\alpha)\,\Delta + e^{-\rho\Delta}\, v(x + f(x,\alpha)\Delta)\right\}.
\]
\medskip
Subtract $v(x)$ from both sides, divide by $\Delta$, and let $\Delta \to 0$ (using $e^{-\rho\Delta} \approx 1 - \rho\Delta$ and $v(x + f\Delta) \approx v(x) + v'(x) f\Delta$):
\[
0 \;=\; \max_\alpha\!\left\{r(x,\alpha) - \rho v(x) + v'(x) f(x,\alpha)\right\}.
\]
\medskip
Rearrange:
\[
\rho v(x) \;=\; \max_\alpha\!\left\{r(x,\alpha) + v'(x) f(x,\alpha)\right\}.
\]
\medskip
This is exactly the HJB.
\end{frame}

\begin{frame}{Stationary HJB and FOCs}
For the optimal growth problem with state $k$, control $c$, return $u(c)$, and transition $\dot k = f(k) - \delta k - c$:
\[
\rho\, v(k) \;=\; \max_c\!\left\{u(c) + v'(k)\,\big(f(k)-\delta k - c\big)\right\}.
\]
\medskip
First-order condition:
\[
u'(c) \;=\; v'(k) \quad\Longrightarrow\quad c^*(k) = (u')^{-1}(v'(k)).
\]
\medskip

Substituting back:
\[
\rho\, v(k) \;=\; u\!\big(c^*(k)\big) + v'(k)\,\big(f(k) - \delta k - c^*(k)\big).
\]
\medskip
A second-order ODE in $v$ to be solved on $k \in [0, \bar k]$ with appropriate boundary conditions.
\end{frame}

\begin{frame}{Poisson income shocks: state space}
\begin{itemize}
\item In the continuous-time Aiyagari, individual income $y$ jumps between two states $\{y_1, y_2\}$ at \emph{Poisson rates} $\lambda_1, \lambda_2$.
\medskip
\item State of an individual: $(a, y_i)$ with $i \in \{1,2\}$.
\medskip
\item Two value functions $v_1(a)$ and $v_2(a)$, one per income state.
\medskip
\item Asset evolves continuously: $\dot a = y_i + r a - c$.
\medskip
\item Income switches at Poisson rate; from state $i$, the agent moves to state $-i$ at rate $\lambda_i$.
\medskip
\item This setup is the natural continuous-time analogue of the discrete-time Aiyagari with Markov income.
\end{itemize}
\end{frame}

\begin{frame}{Deriving the HJB with jump terms}
Over a short interval $\Delta$, with probability $\lambda_i \Delta$ the agent's income jumps from $y_i$ to $y_{-i}$. The Bellman is
\[
v_i(a) \approx \max_c\!\left\{u(c)\Delta + e^{-\rho\Delta}\!\left[(1-\lambda_i\Delta)\, v_i(a + \dot a\Delta) + \lambda_i\Delta\, v_{-i}(a)\right]\right\}.
\]
\medskip

Letting $\Delta \to 0$ and rearranging:
\[
\rho\, v_i(a) \;=\; \max_c\!\left\{u(c) + v_i'(a)\,(y_i + r a - c) + \lambda_i\,\big(v_{-i}(a) - v_i(a)\big)\right\}.
\]
\medskip
\begin{itemize}
\item Three terms: instantaneous utility, drift term (deterministic capital accumulation), jump term (Poisson income switch).
\item This is the HJB with jumps; the system is $i \in \{1,2\}$.
\end{itemize}
\end{frame}

\begin{frame}{Economic intuition of the jump term}
\[
\lambda_i \big(v_{-i}(a) - v_i(a)\big)
\]
\medskip
\begin{itemize}
\item This is the \emph{option value} of switching income state.
\medskip
\item If $v_{-i}(a) > v_i(a)$ (the other income state would yield higher continuation value at the same wealth), the jump term is \emph{positive} --- the agent enjoys an expected gain from the future income switch.
\medskip
\item Conversely, if $v_{-i}(a) < v_i(a)$, the jump term is a cost.
\medskip
\item The Poisson rate $\lambda_i$ controls the strength of this option value.
\medskip
\item Intuition: precautionary saving in continuous time appears partly through this jump term, in addition to the curvature of $u$ along the deterministic drift.
\end{itemize}
\end{frame}

\begin{frame}{Continuous-time Aiyagari: setup}
\textbf{Households.} A continuum of households indexed by current asset $a \in [\underline a, \infty)$ and income $y_i \in \{y_1, y_2\}$ with Poisson switching rates $\lambda_1, \lambda_2$.
\medskip

\textbf{Firm.} Representative firm with $Y = K^\alpha L^{1-\alpha}$ and competitive factor markets:
\[
r = \alpha (K/L)^{\alpha-1} - \delta, \qquad w = (1-\alpha) (K/L)^\alpha.
\]
\medskip

\textbf{Aggregation.} Aggregate capital and labour:
\[
K = \sum_{i=1}^{2} \int a\, g_i(a)\, da, \qquad
L = \sum_{i=1}^{2} y_i \int g_i(a)\, da,
\]
where $g_i(a)$ is the cross-sectional density in income state $i$.
\end{frame}

\begin{frame}{The aggregation loop}
\begin{enumerate}
\item Guess prices $(r, w)$ (or equivalently aggregate $K$).
\medskip
\item Solve the household HJB to get optimal consumption $c_i^*(a)$ and drift $s_i^*(a) = y_i + r a - c_i^*(a)$.
\medskip
\item Solve the KFE for the stationary density $g_i^*(a)$.
\medskip
\item Compute aggregate $K^* = \sum_i \int a\, g_i^*(a)\, da$.
\medskip
\item Update $(r,w)$ from the firm FOCs at $K^*$, and iterate until $K^* = K$.
\end{enumerate}
\medskip

This is structurally identical to the discrete-time Aiyagari nested fixed point, but every step is now an ODE/PDE rather than a discrete-time DP problem.
\end{frame}

\begin{frame}{The Kolmogorov forward (Fokker--Planck) equation}
For the time-dependent density $g_i(a,t)$:
\[
\partial_t g_i(a,t) \;=\; -\partial_a\!\big(s_i(a,t)\, g_i(a,t)\big)
- \lambda_i\, g_i(a,t) + \lambda_{-i}\, g_{-i}(a,t).
\]
\medskip

\begin{itemize}
\item First term: \emph{drift} --- mass moves at speed $s_i(a,t)$ along the $a$-axis.
\item Second term: \emph{Poisson outflow} from state $i$ at rate $\lambda_i$.
\item Third term: \emph{Poisson inflow} from the other state at rate $\lambda_{-i}$.
\medskip
\item Boundary: at the borrowing limit $a = \underline a$, mass cannot leak below; we impose a reflecting boundary.
\medskip
\item Well-posed under standard conditions (Achdou et al. 2022).
\end{itemize}
\end{frame}

\begin{frame}{Reading the KFE: economic intuition}
\begin{itemize}
\item \textbf{Drift term}: the cross-section flows to the right when households are saving ($s_i > 0$) and to the left when they are dissaving.
\medskip
\item \textbf{Jump terms}: at every point in $a$, mass leaks from state $i$ at rate $\lambda_i$ and arrives from $-i$ at rate $\lambda_{-i}$. The two states ``trade'' agents.
\medskip
\item \textbf{Boundary at $\underline a$}: reflecting boundary --- mass that hits the borrowing constraint cannot continue leftward; instead it accumulates at $\underline a$ for as long as $s_i(\underline a) \leq 0$.
\medskip
\item \textbf{Total mass conserved}: $\int g_1 + \int g_2 = 1$ at all times.
\medskip
\item Analogy: water flowing through two interconnected tanks with leaks at known rates.
\end{itemize}
\end{frame}

\begin{frame}{Stationary KFE}
At the stationary equilibrium ($\partial_t = 0$):
\[
0 \;=\; -\partial_a\!\big(s_i(a)\, g_i(a)\big) - \lambda_i\, g_i(a) + \lambda_{-i}\, g_{-i}(a),
\quad i = 1, 2.
\]
\medskip

\begin{itemize}
\item Two coupled ODEs in $g_1, g_2$, plus the normalisation $\int g_1 + \int g_2 = 1$.
\medskip
\item Reflecting boundary at $\underline a$, integrability at $a \to \infty$.
\medskip
\item Standard discretisation gives a sparse linear system $A^\top g = 0$ where $A$ is the discretised HJB infinitesimal generator (next slide).
\end{itemize}
\end{frame}

\begin{frame}{Discretisation and the upwind scheme}
\begin{center}
\includegraphics[width=0.6\textwidth]{vp}
\end{center}
\vspace{-0.3em}
\begin{itemize}
\item Discretise $a$ on a grid $\{a_1, \ldots, a_N\}$. The HJB needs an approximation of $v_i'(a)$.
\item \textbf{Upwind scheme}: pick the forward or backward difference depending on the \emph{sign of the optimal drift}:
\[
v_i'(a_j) \approx \begin{cases}
(v_i(a_{j+1}) - v_i(a_j))/\Delta a & \text{if } s_i(a_j) > 0 \\
(v_i(a_j) - v_i(a_{j-1}))/\Delta a & \text{if } s_i(a_j) < 0
\end{cases}
\]
\item Provably stable (Barles--Souganidis 1991); without upwinding, the scheme is unstable.
\end{itemize}
\end{frame}

\begin{frame}{HJB and KFE: $A$ vs.\ $A^\top$ duality}
After upwind discretisation, the HJB takes the form
\[
\rho\, \mathbf v \;=\; \mathbf u(\mathbf c) + A(\mathbf v)\, \mathbf v,
\]
where $A(\mathbf v) \in \mathbb{R}^{2N\times 2N}$ is the discretised infinitesimal generator (sparse, tridiagonal-plus-jumps).
\medskip

The discretised stationary KFE is
\[
\mathbf 0 \;=\; A(\mathbf v)^{\top}\, \mathbf g.
\]
\medskip
\begin{itemize}
\item The \emph{same} matrix $A$ generates both equations: HJB on the right, KFE on the left (transpose).
\medskip
\item This duality is at the heart of the Achdou et al.\ approach: solving HJB gives you the KFE almost for free.
\end{itemize}
\end{frame}

\begin{frame}{Putting it all together: outer fixed point on $r$}
\begin{enumerate}
\item Guess $r$.
\medskip
\item Solve HJB by iterating
\[
\rho \mathbf v - A(\mathbf v)\mathbf v \;=\; \mathbf u(\mathbf c)
\]
to find $\mathbf v^*, \mathbf c^*, A^*$.
\medskip
\item Solve KFE: find $\mathbf g^*$ in the null space of $A^{*\top}$, normalised to $\mathbf 1^\top \mathbf g^* = 1$.
\medskip
\item Aggregate: $K^* = \sum_j a_j \big(g^*_{1,j} + g^*_{2,j}\big)$.
\medskip
\item Update $r$ from the firm FOC at $K^*$ and iterate until convergence.
\end{enumerate}
\medskip

This is the Achdou--Lasry--Lions--Moll algorithm for continuous-time HA models. It is fast, accurate, and (with deep learning) extends to high-dimensional state spaces.
\end{frame}

\begin{frame}{Transition dynamics: forward--backward system}
For non-stationary problems (e.g.\ MIT shocks, business cycles), HJB and KFE become time-dependent and the system is \emph{forward--backward}:
\[
\begin{cases}
-\partial_t v_i + \rho v_i \;=\; \max_c\,\{\cdots\}, & \text{(backward in time)}\\[0.4em]
\partial_t g_i \;=\; -\partial_a(s_i g_i) - \lambda_i g_i + \lambda_{-i} g_{-i}. & \text{(forward in time)}
\end{cases}
\]
\medskip
\begin{itemize}
\item Boundary in time: $v_i(\cdot, T) = $ terminal condition (often the stationary $v_i^*$); $g_i(\cdot, 0) = $ initial condition.
\medskip
\item Numerical scheme: \emph{shooting} or \emph{Newton's method} on the residuals --- iterate backward HJB given $g$, then forward KFE given $v$, until self-consistent.
\end{itemize}
\end{frame}

\begin{frame}{Master equation: distribution as state}
\begin{itemize}
\item The master equation is a single equation that encodes both HJB and KFE simultaneously, with the \emph{distribution itself} as a state variable.
\medskip
\item Schematically:
\[
\rho\, U(x, m) \;=\; H\!\big(x, \nabla_x U(x,m), m\big) + \int \nabla_m U(x,m)(y)\, \mathcal L^* m(y)\, dy,
\]
where $U$ is the value function defined on the product space (individual state, distribution) and $\mathcal L^*$ is the dual generator on the distribution.
\medskip
\item This is an infinite-dimensional PDE on the Wasserstein space of probability measures.
\medskip
\item Practical use: linearisation around the stationary equilibrium gives a tractable system for solving HA models with aggregate shocks (Bilal 2023; Alvarez--Lippi 2022).
\end{itemize}
\end{frame}

\begin{frame}{Deep learning for MFG}
\begin{itemize}
\item Classical numerical schemes (upwind, finite differences) suffer in high-dimensional state spaces.
\medskip
\item \textbf{Deep neural networks for MFG}:
  \begin{itemize}
  \item Parameterise $v(x, m)$ and $m(x)$ by NNs.
  \item Train by minimising the residual of the HJB and the KFE jointly.
  \item Use Monte Carlo sampling instead of grids.
  \end{itemize}
\medskip
\item Recent progress (Carmona--Lauri\`ere 2021, Lauri\`ere et al.\ 2022) shows that deep MFG solvers can handle 10--20 dimensional state spaces that are completely out of reach for grid methods.
\medskip
\item These methods unify the discrete-time RL approaches (DeepHAM, DEQN) with the continuous-time PDE perspective.
\end{itemize}
\end{frame}

\begin{frame}{MFG $\leftrightarrow$ HA models: the link}
\begin{itemize}
\item Discrete-time Aiyagari is the natural \emph{discretisation} of continuous-time MFG with idiosyncratic Poisson income.
\medskip
\item Krusell--Smith with aggregate uncertainty corresponds to MFGs with a \emph{common noise} (the aggregate shock).
\medskip
\item In both formulations:
  \begin{itemize}
  \item Optimal individual behaviour depends on the cross-sectional distribution.
  \item The cross-sectional distribution evolves under optimal behaviour.
  \item Equilibrium is a fixed point of this circular dependence.
  \end{itemize}
\medskip
\item MFG provides the cleanest mathematical formalism, while RL/DL provide the computational engine.
\medskip
\item The two approaches are now converging: deep MFG solvers borrow heavily from RL; modern deep RL borrows from continuous-time control theory.
\end{itemize}
\end{frame}

%=============================================================================
% PART XI: LAB PREVIEW & WRAP-UP
%=============================================================================

\section{Lab Preview and Wrap-up}

\begin{frame}{Lab06 preview: Krusell--Smith via Deep Equilibrium Nets}
\begin{itemize}
\item In the 1-hour lab you will train a Deep Equilibrium Net to solve the Krusell--Smith (1998) economy in PyTorch.
\medskip
\item Notebook: \texttt{Lab06\_KS\_HA.ipynb} (in this folder).
\medskip
\item By the end of the lab you will:
  \begin{itemize}
  \item Initialise the model from a steady state computed via the sequence-Jacobian (SSJ) toolbox.
  \item Build a PyTorch policy network for the savings rate.
  \item Run the cloud-of-economies training loop.
  \item Visualise the trained policy and the loss curves.
  \item Diagnose convergence and accuracy.
  \end{itemize}
\medskip
\item Hardware: a single GPU (Colab T4 / V100) is comfortable; CPU also works for short runs.
\end{itemize}
\end{frame}

\begin{frame}{Lab06 step-by-step}
\textbf{Part 1.} Initialise via sequence-Jacobian:
\begin{itemize}
\item Solve the steady state to get asset and employment grids, transition matrices, $K_{ss}, \mu_{ss}$.
\end{itemize}
\medskip

\textbf{Part 2.} PyTorch implementation:
\begin{itemize}
\item Define the policy network (3 hidden layers, Mish, sigmoid output).
\item Encode the budget constraint inside the network.
\item Wrap the model physics (transition dynamics) as a class.
\end{itemize}
\medskip

\textbf{Part 3.} The training loop (cloud simulation):
\begin{itemize}
\item Simulate $N$ parallel economies.
\item Compute the Euler-residual loss.
\item Backprop and Adam update.
\item Transport the distribution forward (no grad).
\end{itemize}
\medskip

\textbf{Part 4.} Analysis:
\begin{itemize}
\item Plot the loss history (log scale).
\item Plot policy functions; check Euler errors.
\end{itemize}
\end{frame}

\begin{frame}{Optional follow-up: Lab12\_OLG.ipynb}
\begin{itemize}
\item Once you finish the KS lab, you may want to explore \texttt{Notebooks/Lab12\_OLG.ipynb}.
\medskip
\item This notebook implements OLG with aggregate uncertainty using JAX and stabilising homotopies (Brumm--Kubler--Scheidegger 2023).
\medskip
\item Highlights:
  \begin{itemize}
  \item Two-stage homotopy: first solve a capital-only model, then introduce bonds.
  \item JAX \texttt{jit} compilation for fast training-simulation episodes.
  \item Detailed pedagogical commentary inside the notebook.
  \end{itemize}
\medskip
\item Useful preparation for research projects involving life-cycle, fiscal policy, or asset-pricing applications.
\end{itemize}
\end{frame}

\begin{frame}{Recap of Lecture 6}
\begin{itemize}
\item \textbf{Aiyagari--Bewley--Huggett}: incomplete-markets HA model and stationary equilibrium; existence/uniqueness via Feller, mixing, monotonicity.
\medskip
\item \textbf{Classical computation}: nested fixed-point on $r$; VFI on the household, eigenvector for the distribution, bisection on the price.
\medskip
\item \textbf{Why ML?} Curse of dimensionality once we add aggregate shocks, multiple assets, or life-cycle structure.
\medskip
\item \textbf{Deep dynamic programming}: policy/value networks, Euler residual losses, actor--critic.
\medskip
\item \textbf{DeepHAM and DEQN}: two leading deep-learning solutions for KS.
\medskip
\item \textbf{OLG with aggregate uncertainty}: homotopy-stabilised JAX training.
\medskip
\item \textbf{Mean-Field Games}: continuous-time formulation, HJB+KFE, deep MFG solvers as the unifying frontier.
\medskip
\item \textbf{Next lecture}: Large Language Models --- a different but equally fast-moving deep-learning frontier.
\end{itemize}
\end{frame}

\end{document}
